\section{总结与延伸}

\subsection{Context Engineering 的全景图}

本期完成了 Codex Context Engineering 的全景分析：

\begin{itemize}
    \item \textbf{Skills}：出现在 User Message 的 content 中，采用渐进式加载
    \item \textbf{Tools}：出现在 API 的 tools 字段中，包括内置 Tools 和 MCP 动态 Tools
    \item \textbf{System 指令}：独立的 instructions 字段
    \item \textbf{消息历史}：多轮交互记录，可通过 compact 压缩
    \item \textbf{Memory}：持久化的记忆系统
\end{itemize}

\begin{importantbox}{三个关键认知}
\begin{enumerate}
    \item 不管是 Skills、MCP 还是预设 Tools，如果不暴露在 Context 中，模型就无法感知和调用
    \item Skills 和 Tools 在技术上正交独立——Skills 在消息内容中，Tools 在工具字段中
    \item 所有 Modern Agent 的核心工作都是 Context Engineering——精心管理送给模型的上下文
\end{enumerate}
\end{importantbox}

\subsection{拓展阅读}

\begin{itemize}
    \item OpenAI Codex 官方文档中的 Agent Loop 设计
    \item Claude Code 的 CLAUDE.md 和 Skills 机制
    \item MCP（Model Context Protocol）规范
    \item OpenAI Response API 文档
\end{itemize}

%% --- Fin del contenido --- %%

\end{document}
